% ============================================================
% Section 187: 铜自由电子气——费米能、漂移速度与能态密度
% ============================================================
\section{固体理论：铜自由电子气——费米能、漂移速度与能态密度}
\label{sec:cu-free-electron}

\begin{modelbox}[题目：铜的费米能与 Drude 输运参量]
铜的密度为 $9.9\times10^3\,\text{kg/m}^3$，电阻率为 $1.56\times10^{-8}\,\Omega\cdot\text{cm}$，原子量为 $63.5$，原子密度为 $8.4\times10^{22}\,\text{cm}^{-3}$。已知 $\hbar=1.055\times10^{-34}\,\text{J}\cdot\text{s}$，电子质量 $m=9.1\times10^{-31}\,\text{kg}$，电子电量 $e=1.6\times10^{-19}\,\text{C}$。假设一个铜原子有一个价电子：
\begin{enumerate}
    \item 求绝对零度下铜自由电子的费米能 $E_F$。若外加 $1.0\,\text{V/cm}$ 直流电场，求自由电子平均漂移速度、弛豫时间和平均自由程；
    \item 分别导出一维、二维、三维铜金属中自由电子气的能态密度。
\end{enumerate}
\end{modelbox}

\begin{tipbox}[考点快速分析]
\begin{itemize}
    \item \textbf{费米能公式}：$E_F=\frac{\hbar^2}{2m}(3\pi^2 n)^{2/3}$
    \item \textbf{Drude 弛豫时间}：$\rho=\frac{m}{ne^2\tau}\;\Rightarrow\;\tau=\frac{m}{ne^2\rho}$
    \item \textbf{漂移速度}：$v_d=\frac{eE\tau}{m}=\frac{E}{ne\rho}$（由 $j=ne v_d=\sigma E$）
    \item \textbf{平均自由程}：$l=v_F\tau$，费米速度 $v_F=\sqrt{2E_F/m}$
    \item \textbf{能态密度}：$g(E)=\frac{\dd N}{\dd E}$，关键是把 $k$ 空间体积元换算到能量
\end{itemize}
\end{tipbox}

\begin{derivationbox}[标准解题过程]
\textbf{(1) 费米能、弛豫时间、漂移速度与平均自由程}

电子浓度 $n=8.4\times10^{22}\,\text{cm}^{-3}=8.4\times10^{28}\,\text{m}^{-3}$。

\textbf{费米能}：
\begin{equation}
3\pi^2 n=3\pi^2(8.4\times10^{28})\approx2.49\times10^{30}\,\text{m}^{-3},
\quad (3\pi^2 n)^{2/3}\approx1.84\times10^{20}\,\text{m}^{-2},
\end{equation}
\begin{equation}
E_F=\frac{(1.055\times10^{-34})^2}{2\times9.1\times10^{-31}}\times1.84\times10^{20}
\approx\frac{1.113\times10^{-68}}{1.82\times10^{-30}}\times1.84\times10^{20}
\approx1.12\times10^{-18}\,\text{J}\approx\boxed{7.0\,\text{eV}.}
\end{equation}

\textbf{弛豫时间}（Drude 公式）：
\begin{equation}
\tau=\frac{m}{ne^2\rho}
=\frac{9.1\times10^{-31}}{8.4\times10^{28}\times(1.6\times10^{-19})^2\times1.56\times10^{-10}}
\approx\frac{9.1\times10^{-31}}{3.355\times10^{-19}}
\approx\boxed{2.7\times10^{-12}\,\text{s}.}
\end{equation}
（注：实际铜的弛豫时间约 $10^{-14}\sim10^{-13}\,\text{s}$ 量级，本题因给定数据计算结果偏大。）

\textbf{漂移速度}：$E=1.0\,\text{V/cm}=100\,\text{V/m}$，
\begin{equation}
v_d=\frac{eE\tau}{m}
=\frac{1.6\times10^{-19}\times100\times2.7\times10^{-12}}{9.1\times10^{-31}}
\approx\frac{4.32\times10^{-29}}{9.1\times10^{-31}}
\approx\boxed{47.5\,\text{m/s}.}
\end{equation}

\textbf{费米速度}：
\begin{equation}
v_F=\sqrt{\frac{2E_F}{m}}
=\sqrt{\frac{2\times1.12\times10^{-18}}{9.1\times10^{-31}}}
\approx\sqrt{2.46\times10^{12}}
\approx1.57\times10^6\,\text{m/s}.
\end{equation}

\textbf{平均自由程}：
\begin{equation}
l=v_F\tau\approx1.57\times10^6\times2.7\times10^{-12}
\approx\boxed{4.2\times10^{-6}\,\text{m}=4.2\,\mu\text{m}.}
\end{equation}

\textbf{(2) 各维度能态密度}

自由电子色散 $E=\frac{\hbar^2k^2}{2m}$，$k=\sqrt{2mE}/\hbar$，$\dd E=\frac{\hbar^2k}{m}\dd k$。

\begin{itemize}
    \item \textbf{三维}：$k$ 空间体积元 $2\times\frac{V}{(2\pi)^3}4\pi k^2\dd k$（因子2为自旋），
    \begin{equation}
    g_{3D}(E)=\frac{V}{2\pi^2}\Bigl(\frac{2m}{\hbar^2}\Bigr)^{3/2}\sqrt{E}.
    \end{equation}
    \item \textbf{二维}：$k$ 空间面积元 $2\times\frac{A}{(2\pi)^2}2\pi k\dd k$，
    \begin{equation}
    g_{2D}(E)=\frac{Am}{\pi\hbar^2}=\text{常数}.
    \end{equation}
    \item \textbf{一维}：$k$ 空间线元 $2\times\frac{L}{2\pi}\dd k$（两支），
    \begin{equation}
    g_{1D}(E)=\frac{L}{\pi}\sqrt{\frac{2m}{\hbar^2}}\frac{1}{\sqrt{E}}.
    \end{equation}
\end{itemize}
\end{derivationbox}

\begin{insightbox}[物理图像]
\begin{itemize}
    \item 费米能 $E_F\sim7\,\text{eV}$ 是室温 $k_BT\sim0.025\,\text{eV}$ 的约 $300$ 倍，说明铜中自由电子是强简并的费米气体，量子效应占主导。
    \item 漂移速度 $v_d\sim50\,\text{m/s}$ 远小于费米速度 $v_F\sim10^6\,\text{m/s}$，电场仅使费米球整体微小偏移，绝大部分电子仍在费米面附近做高速无规运动。
    \item 平均自由程 $l\sim\mu\text{m}$ 在低温纯净样品中可达 mm 量级，这是金属中电子波动性的直接体现。
\end{itemize}
\end{insightbox}

\begin{formulabox}[考场速记]
\textbf{Drude 公式链}：$\sigma=\frac{ne^2\tau}{m}=ne\mu=\frac{1}{\rho}$，迁移率 $\mu=\frac{e\tau}{m}=\frac{v_d}{E}$。

\textbf{三维费米参量关系}：
\begin{center}
\begin{tabular}{lll}
\toprule
\textbf{量} & \textbf{公式} & \textbf{数量级（铜）} \\
\midrule
$E_F$ & $\frac{\hbar^2}{2m}(3\pi^2 n)^{2/3}$ & $7\,\text{eV}$ \\
$k_F$ & $(3\pi^2 n)^{1/3}$ & $1.36\times10^{10}\,\text{m}^{-1}$ \\
$v_F$ & $\hbar k_F/m$ & $1.57\times10^6\,\text{m/s}$ \\
$\lambda_F$ & $2\pi/k_F$ & $0.46\,\text{nm}$ \\
\bottomrule
\end{tabular}
\end{center}
\end{formulabox}
